\section{实验与精度分析}
\label{sec:exp}

本节围绕三条主线展开：一是验证系统的重复精度（同一位置多次测量的一致性）是否
达标，二是通过独立高精度参考验证测量的绝对准确度，三是分析测量精度受哪些因素
影响并定位主导误差来源。需要说明的是，本节部分定量结果仍在补充实验中，相应数值
以占位形式标注，不作为最终结论。

\subsection{实验装置与数据采集}

样机由前、后两个相机模块沿设备 $X$ 方向并排组成，两模块中心间距约
$80\,\mathrm{mm}$。每个模块含一片 $45^\circ$ 反射镜（实测夹角 $45.22^\circ$、
厚约 $1.05\,\mathrm{mm}$、两面平行度优于 $0.06^\circ$）与一块紧贴取景框的毛
玻璃接收屏（取景框与毛玻璃叠层总厚约 $5\,\mathrm{mm}$）。相机分辨率为
$2592\times1944$。测杆固定在底板上方，靶点位于杆末端。设备各面的绝对位置与
反射镜、测杆的几何参数由影像仪与关节臂配合标定，作为设备坐标系几何真值的对照
基准。采集时对同一目标位置连续突发拍摄 $N$ 帧、取圆心均值以抑制随机噪声，并
按“杆长—距离”分组多次重复，用于统计重复精度。真值参考由关节臂扫描给出。

\subsection{重复精度}

重复精度以每组重复测量的组内标准差 $\hat{\sigma}_H$ 衡量。当前基线下，全体
测量组的平均组内标准差约为 \todo{重复精度 $\hat{\sigma}_H$ 终值}，与工程目标
$\pm15$--$50\,\mu\mathrm{m}$ 之间的关系为 \todo{是否达标的结论}。在优化过程中
发现，圆心提取策略对重复精度具有决定性影响：激光光斑并非小圆点，而是半径约
$120\,\mathrm{px}$、中心大面积饱和且带有不对称暗弱拖尾的大斑；早期采用的小
窗口、低阈值强度加权重心会把拖尾纳入重心而被拉偏，改为覆盖整斑的大窗口配合高
阈值后，平均组内标准差从约 $216\,\mu\mathrm{m}$ 降至约 $106\,\mu\mathrm{m}$，
接近腰斩。这一改进已合入测量流程。相关分组重复精度的完整统计见
表~\ref{tab:repeat}。

\begin{table}[t]
\centering
\caption{分组重复精度统计（组内标准差 $\hat{\sigma}_H$）}
\label{tab:repeat}
\begin{tabular}{@{}lccc@{}}
\toprule
分组（杆长／距离） & 重复次数 & $\hat{\sigma}_H/\mu\mathrm{m}$ & 备注 \\
\midrule
\todo{组1} & \todo{$n$} & \todo{$\sigma$} & \\
\todo{组2} & \todo{$n$} & \todo{$\sigma$} & \\
平均       & --          & \todo{均值}     & \\
\bottomrule
\end{tabular}
\end{table}

\subsection{绝对值验证}

重复精度只反映同一位置多次测量的一致性，即测量的稳定性，却不能说明测得的距离
是否接近真值——一个系统可能重复性很好却存在稳定的系统偏差。因此还需对测量的
绝对准确度单独验证。由于“空间任意点到激光基准轴距离”这一测量量此前尚无成熟的
标准量具或公认的参考方法，我们无法直接购置一台已校准的仪器作真值来源，只能自行
搭建一套独立的高精度参考实验。

我们计划在实验室环境下借助激光跟踪仪构建参考基准。如第~\ref{sec:related} 节所
分析，激光跟踪仪的三维定位精度显著优于本文目标，虽因需逐点安放合作反射器而不适
用于现场高密度对中，却恰好适合在实验室对少量标定点提供独立、可溯源的高精度坐标
参考。具体方案是：用激光跟踪仪配合合作反射器测量激光基准轴上若干点以独立拟合出
基准轴，再测量各目标点坐标并按点到直线公式算出参考距离 $H_{\mathrm{ref}}$；将
DCPAM 在同一批目标点上测得的 $H_{\mathrm{DCPAM}}$ 与 $H_{\mathrm{ref}}$ 逐点
比较，以偏差 $H_{\mathrm{DCPAM}}-H_{\mathrm{ref}}$ 的均值刻画系统性误差（trueness）、
以其标准差与重复精度相互印证，从而给出 DCPAM 的绝对准确度评估。

该实验尚在进行中，结果留待补充：绝对偏差均值为 \todo{系统偏差均值}，偏差标准差
为 \todo{偏差标准差}，与工程目标 $\pm15$--$50\,\mu\mathrm{m}$ 的关系为
\todo{绝对准确度是否达标的结论}。

\subsection{误差来源与灵敏度分析}

测量误差可归为测量端与几何端两类。测量端的主导来源是圆心提取的像素噪声：即便
在亚像素级（跨图圆心标准差约 $0.3$--$0.7\,\mathrm{px}$），其经由整条链路传播后
仍构成 $\hat{\sigma}_H$ 的主要成分。这一噪声进一步被光路几何放大——分析表明，
组内标准差与虚像点在 $Z$ 方向的离散强相关（相关系数约 $0.944$），像素扰动经
约 $740\,\mathrm{mm}$ 的杠杆臂放大后显著抬升了 $H$ 的波动，因此 $Z$ 方向的杠杆
敏感度是需要在光路与靶点几何设计上重点抑制的因素。此外，激光器自身的指向抖动、
相机在采集期的机械抖动，以及曝光不一致带来的采集期漂移，都会叠加到圆心的帧间
波动上，其量级需通过专门实验分别标定，此处以占位标注：激光抖动
\todo{激光抖动贡献}、相机抖动 \todo{相机抖动贡献}、PnP 定位精度
\todo{PnP 复现精度}。

几何端的敏感度则相对次要。围绕反射面平移与设备几何参数所做的偏移扫描与自标定
实验表明，当前设备几何已接近物理真值：把反射镜夹角、模块间距、装配旋转、测杆
杆根与杆长等参数放开做以组内方差为目标的自标定，各参数仍稳定落在物理真值附近
（如反射镜夹角约 $45^\circ$、模块间距约 $80\,\mathrm{mm}$、装配旋转小于
$0.01^\circ$），而平均组内标准差仅改善约 $5\%$。这从反面印证了两点：其一，
几何标定已到位、无需再调；其二，$\hat{\sigma}_H$ 的主导来源并非几何参数，而是
测量端的像素噪声。因此，降低方差的主战场在测量端而非几何参数。

\subsection{误差传播与蒙特卡洛验证}

基于第~\ref{sec:model} 节的端到端仿射形式，把 $H$ 对四个像素坐标的一阶泰勒
展开写为
\begin{equation}
\label{eq:propagation}
\sigma_H^2 \approx \sum_{q\in\{u_1,v_1,u_2,v_2\}}
\left(\frac{\partial H}{\partial q}\right)^2 \sigma_q^2,
\end{equation}
并定义各像素坐标的方差贡献占比
\begin{equation}
\label{eq:eta}
\eta_q = \frac{(\partial H/\partial q)^2\,\sigma_q^2}{\sigma_H^2}\times 100\%.
\end{equation}
为验证解析式，对每个像素坐标在 $\pm1\,\mu\mathrm{m}$ 范围内独立扰动做
$N=10^4$ 次蒙特卡洛采样，比较经验标准差 $\hat{\sigma}_H$ 与解析值 $\sigma_H$。
二者的一致性结果为 \todo{解析 $\sigma_H$ 与 MC $\hat{\sigma}_H$ 对照值}，
各像素坐标的灵敏度占比 $\eta_q$ 为 \todo{$\eta_q$ 分布}，与前述“$Z$ 向杠杆
放大像素噪声”的判断相互印证。

\subsection{讨论}

综合上述结果，当前系统的重复精度与 $\pm15$--$50\,\mu\mathrm{m}$ 的工程目标之间
\todo{差距量化}，其瓶颈明确位于测量端的像素噪声及其经 $Z$ 向杠杆的放大，而非
几何标定。据此，后续改进应从三方面着手：进一步提升圆心提取精度或改用更高信噪比
的相机与激光源，以直接压低像素噪声；通过光路与靶点几何设计降低 $Z$ 方向的杠杆
敏感度；以及改善采集条件（统一曝光、抑制采集期漂移与机械抖动）。这些方向所对应
的定量收益将在补充实验中给出。
